\documentclass[10pt,a4paper]{amsart}
\usepackage{amsmath,amssymb,amsthm,graphicx}
% Retained minimal preamble for the Stacks Project chapter more-morphisms.tex (More on Morphisms).
% Verbatim excerpt of the chapter preamble: Xy-pic setup (native preamble lines 14-18) and the native
% theorem declarations with the mathematical macros (native preamble lines 154-217). Copied mechanically;
% nothing added, removed or redefined. The wrapper class and the amsmath/amssymb/amsthm/graphicx base
% packages are supplied by the assembler, so the class/externaldocument part of the native preamble is omitted.
% For commutative diagrams we use Xy-pic
\usepackage[all]{xy}

% We use 2cell for 2-commutative diagrams.
\xyoption{2cell}
% Theorem environments.
%
\theoremstyle{plain}
\newtheorem{theorem}[subsection]{Theorem}
\newtheorem{proposition}[subsection]{Proposition}
\newtheorem{lemma}[subsection]{Lemma}

\theoremstyle{definition}
\newtheorem{definition}[subsection]{Definition}
\newtheorem{example}[subsection]{Example}
\newtheorem{exercise}[subsection]{Exercise}
\newtheorem{situation}[subsection]{Situation}

\theoremstyle{remark}
\newtheorem{remark}[subsection]{Remark}
\newtheorem{remarks}[subsection]{Remarks}

\numberwithin{equation}{subsection}

% Macros
%
\def\lim{\mathop{\mathrm{lim}}\nolimits}
\def\colim{\mathop{\mathrm{colim}}\nolimits}
\def\Spec{\mathop{\mathrm{Spec}}}
\def\Hom{\mathop{\mathrm{Hom}}\nolimits}
\def\Ext{\mathop{\mathrm{Ext}}\nolimits}
\def\SheafHom{\mathop{\mathcal{H}\!\mathit{om}}\nolimits}
\def\SheafExt{\mathop{\mathcal{E}\!\mathit{xt}}\nolimits}
\def\Sch{\mathit{Sch}}
\def\Mor{\mathop{\mathrm{Mor}}\nolimits}
\def\Ob{\mathop{\mathrm{Ob}}\nolimits}
\def\Sh{\mathop{\mathit{Sh}}\nolimits}
\def\NL{\mathop{N\!L}\nolimits}
\def\CH{\mathop{\mathrm{CH}}\nolimits}
\def\proetale{{pro\text{-}\acute{e}tale}}
\def\etale{{\acute{e}tale}}
\def\QCoh{\mathit{QCoh}}
\def\Ker{\mathop{\mathrm{Ker}}}
\def\Im{\mathop{\mathrm{Im}}}
\def\Coker{\mathop{\mathrm{Coker}}}
\def\Coim{\mathop{\mathrm{Coim}}}

% Boxtimes
%
\DeclareMathSymbol{\boxtimes}{\mathbin}{AMSa}{"02}

%
% Macros for moduli stacks/spaces
%
\def\QCohstack{\mathcal{QC}\!\mathit{oh}}
\def\Cohstack{\mathcal{C}\!\mathit{oh}}
\def\Spacesstack{\mathcal{S}\!\mathit{paces}}
\def\Quotfunctor{\mathrm{Quot}}
\def\Hilbfunctor{\mathrm{Hilb}}
\def\Curvesstack{\mathcal{C}\!\mathit{urves}}
\def\Polarizedstack{\mathcal{P}\!\mathit{olarized}}
\def\Complexesstack{\mathcal{C}\!\mathit{omplexes}}
% \Pic is the operator that assigns to X its picard group, usage \Pic(X)
% \Picardstack_{X/B} denotes the Picard stack of X over B
% \Picardfunctor_{X/B} denotes the Picard functor of X over B
\def\Pic{\mathop{\mathrm{Pic}}\nolimits}
\def\Picardstack{\mathcal{P}\!\mathit{ic}}
\def\Picardfunctor{\mathrm{Pic}}
\def\Deformationcategory{\mathcal{D}\!\mathit{ef}}
\usepackage{cleveref}
\crefname{originalreference}{original reference}{original references}
\Crefname{originalreference}{Original reference}{Original references}
\makeatletter
\newlabel{morphisms-definition-family-ample-invertible-modules}{{0FXR}{source}{}{sourcepointer-1}{}}
\newlabel{morphisms-definition-family-ample-invertible-modules@cref}{{[originalreference][0][]0FXR}{[1][0][]source}{}{}{}}
\makeatother
\crefname{originalreference}{original reference}{original references}
\Crefname{originalreference}{Original reference}{Original references}
\makeatletter
\newlabel{morphisms-definition-relatively-ample}{{01VH}{source}{}{sourcepointer-2}{}}
\newlabel{morphisms-definition-relatively-ample@cref}{{[originalreference][0][]01VH}{[1][0][]source}{}{}{}}
\makeatother
\crefname{originalreference}{original reference}{original references}
\Crefname{originalreference}{Original reference}{Original references}
\makeatletter
\newlabel{properties-lemma-invert-s-sections}{{01PW}{source}{}{sourcepointer-3}{}}
\newlabel{properties-lemma-invert-s-sections@cref}{{[originalreference][0][]01PW}{[1][0][]source}{}{}{}}
\makeatother
\title{An ample family on X from an ample family on the base and an f-ample invertible module}
\author{The Stacks Project Authors}
\date{stacks-project:more-morphisms.tex:local-snapshot-20261011; GNU FDL 1.2 or later}
\begin{document}
\maketitle
\noindent\textit{Editorial note.} Counted span is the author's own proof environment, source lines 23312-23335, immediately following the statement at 23302-23310.\par\medskip
\noindent\textit{Editorial note.} One cross-chapter pointer (Properties) is recorded in assembly.external\_references; the rest of the inputs are definitional citations inside the chapter.\par\medskip
\noindent\textit{Editorial note.} Related card S9-morph-get-ample-family (L26019) treats the different statement on Cartier divisors; the two share no span.\par\medskip
\noindent\textit{Allowed standard prerequisites (editorial).}\begin{itemize}
\item Definition of an ample family of invertible modules. Morphisms, Definition morphisms-definition-family-ample-invertible-modules, cited at proof lines 23314-23319. Fixes what has to be produced on X: enough sections whose nonvanishing loci are affine and cover.
\item Relative ampleness. Morphisms, Definition morphisms-definition-relatively-ample, cited at proof lines 23312-23314. f-ample gives quasi-compactness of f and the ample-over-preimage behaviour of L used in the local step.
\item Extending sections after inverting a function. Properties, Lemma properties-lemma-invert-s-sections, cited at proof lines 23325-23327. Extends the local section s of L over X\_\{f\textasciicircum{}*t\} to a global section s' of f\textasciicircum{}*N\textasciicircum{}\{tensor n\} tensor L after multiplying by (f\textasciicircum{}*t)\textasciicircum{}n.
\end{itemize}\medskip
\ifdefined\hypertarget\hypertarget{sourcepointer-1}{}\fi
\noindent\textit{Source reference.} Original reference 0FXR: focused-seeds/geometry/S9-morph-ample-family-ample-relative/ref-tags.txt, source lines 5--5. Navigation pointer only: the counted proof cites this statement by \textbackslash{}ref as a standing input; the deduction the counted proof has to make is carried out inside the counted span, and no part of that deduction is delegated to this pointer.\par\medskip
\ifdefined\hypertarget\hypertarget{sourcepointer-2}{}\fi
\noindent\textit{Source reference.} Original reference 01VH: focused-seeds/geometry/S9-morph-ample-family-ample-relative/ref-tags.txt, source lines 6--6. Navigation pointer only: the counted proof cites this statement by \textbackslash{}ref as a standing input; the deduction the counted proof has to make is carried out inside the counted span, and no part of that deduction is delegated to this pointer.\par\medskip
\ifdefined\hypertarget\hypertarget{sourcepointer-3}{}\fi
\noindent\textit{Source reference.} Original reference 01PW: focused-seeds/geometry/S9-morph-ample-family-ample-relative/ref-tags.txt, source lines 7--7. Navigation pointer only: the counted proof cites this statement by \textbackslash{}ref as a standing input; the deduction the counted proof has to make is carried out inside the counted span, and no part of that deduction is delegated to this pointer.\par\medskip
\begin{lemma}
\label{lemma-ample-family-ample-relative}
Let $f : X \to Y$ be a morphism of schemes. Assume
\begin{enumerate}
\item $Y$ has an ample family of invertible modules,
\item there exists an $f$-ample invertible module on $X$.
\end{enumerate}
Then $X$ has an ample family of invertible modules.
\end{lemma}

\begin{proof}
Let $\mathcal{L}$ be an $f$-ample invertible module on $X$.
This in particular implies that $f$ is quasi-compact, see
Morphisms, Definition \ref{morphisms-definition-relatively-ample}.
Since $Y$ is quasi-compact by Morphisms, Definition
\ref{morphisms-definition-family-ample-invertible-modules}
we see that $X$ is quasi-compact (and hence $X$ itself satisfies the
first condition of  Morphisms, Definition
\ref{morphisms-definition-family-ample-invertible-modules}).
Let $x \in X$ with image $y \in Y$. By assumption (2) we can
find an invertible $\mathcal{O}_Y$-module $\mathcal{N}$ and a section
$t \in \Gamma(Y, \mathcal{N})$ such that the locus $Y_t$ where $t$
does not vanish is affine. Then $\mathcal{L}$ is ample over
$f^{-1}(Y_t) = X_{f^*t}$ and hence we can find a section
$s \in \Gamma(X_{f^*t}, \mathcal{L})$ such that $(X_{f^*t})_s$ is affine
and contains $x$. By
Properties, Lemma \ref{properties-lemma-invert-s-sections}
for some $n \geq 0$ the product $(f^*t)^n s$ extends to a section
$s' \in \Gamma(X, f^*\mathcal{N}^{\otimes n} \otimes \mathcal{L})$.
Then finally the section $s'' = f^* ts'$ of
$f^*\mathcal{N}^{\otimes n + 1} \otimes \mathcal{L}$
vanishes at every point of $X \setminus X_{f^*t}$ hence
we see that $X_{s''} = (X_{f^*t})_s$ is affine as desired.
\end{proof}

\end{document}
